\documentclass[11pt,letterpaper]{article}
\usepackage[utf8]{inputenc}
\usepackage{geometry}
\usepackage{amsmath}
\usepackage{setspace}
\usepackage{url}
\usepackage{enumitem}
\usepackage{hyperref}

\geometry{left=1.0in,right=1.0in,top=1.0in,bottom=1.0in}
\onehalfspacing

\setlength{\parindent}{0pt}
\setlength{\parskip}{1.0em plus 0.2em minus 0.2em}

\begin{document}

\thispagestyle{empty}

\begin{center}
    \textbf{\large IN THE UNITED STATES PATENT AND TRADEMARK OFFICE}
\end{center}

\vspace{0.3cm}
\noindent\begin{tabular}{p{2.3in}p{4.0in}}
\textbf{Applicant / Inventor:} & Reinaldo Javier Aguiar Marcano \\
\textbf{Application Number:} & 19/562,829 \\
\textbf{Filing Date:} & March 11, 2026 \\
\textbf{Confirmation Number:} & 5894 \\
\textbf{Art Unit / Examiner:} & TC 2100 / ANDREA BENNETT \\
\textbf{Office of Petitions:} & Cliff Congo, Attorney Advisor \\
\textbf{Customer Number:} & 226458 \\
\textbf{Entity Status:} & SMALL ENTITY (37 CFR \S{} 1.27) \\
\textbf{Title of Invention:} & METHOD FOR HIGH-RADIX COMPUTING VIA OPTICAL BUS \\
\end{tabular}

\vspace{0.4cm}
\begin{center}
    \textbf{\large APPLICANT RESPONSE TO NOTICE OF INCOMPLETE REPLY \\
    AND NOTICE TO FILE MISSING PARTS OF NONPROVISIONAL APPLICATION} \\
    \textbf{\normalsize (SUBMISSION OF TWO-MONTH EXTENSION OF TIME PETITION, \\
    CONDITIONAL SUSPENSION, AND RESPECTFUL REQUEST FOR DEPOSIT ALLOCATION)}
\end{center}

\noindent\textbf{Mail Stop Missing Parts} \\
Commissioner for Patents \\
P.O. Box 1450 \\
Alexandria, VA 22313-1450

\vspace{0.2cm}
Sir:

In responsive reply to the \textit{Notice of Incomplete Reply} mailed August 27, 2026, the \textit{Notice to File Missing Parts} mailed July 9, 2026, and the \textit{Decision on Petition} (PETDEC) mailed September 29, 2026, the pro se Applicant, Reinaldo Javier Aguiar Marcano, respectfully submits this comprehensive formalities response and fee reconciliation.

\section*{I. STATEMENT OF REAFFIRMED ENTITY STATUS (37 CFR \S{} 1.27)}

Applicant hereby formally reaffirms entitlement to \textbf{Small Entity status} under 37 CFR \S{} 1.27 and 35 U.S.C. \S{} 123. Applicant's verified Statement of Small Entity Status was filed on August 2, 2026, and officially recognized and entered on the docket by the Office. All statutory fees are computed strictly under the Small Entity schedule.

\section*{II. PROCEDURAL POSTURE \& CURED FORMALITIES STATUS}

Applicant respectfully notes the following procedural milestones established in the Image File Wrapper:
\begin{enumerate}[leftmargin=*]
    \item \textbf{Cured Technical Informalities:} In response to the July 9, 2026 Notice to File Missing Parts, Applicant submitted replacement claims commencing on a separate sheet and a replacement abstract commencing on a separate sheet on August 6, 2026. The Office officially confirmed entry of these compliant technical papers in the August 27, 2026 Notice of Incomplete Reply.
    \item \textbf{Inventor's Declaration:} An executed Inventor's Oath or Declaration under 37 CFR \S{} 1.63 was filed on July 21, 2026, and acknowledged in the record.
    \item \textbf{Sole Outstanding Item:} The August 27, 2026 Notice of Incomplete Reply establishes that the sole remaining items required to advance the application are the statutory pre-examination fee balance and the requisite extension of time under 37 CFR \S{} 1.136(a).
\end{enumerate}

\section*{III. STATUTORY FEE RECONCILIATION, REQUEST FOR DEPOSIT ALLOCATION (37 CFR \S{} 1.25(b)), AND MACRO PORTFOLIO EQUITY}

\subsection*{A. Macro Portfolio Fee Equity Context}
Applicant respectfully directs the Office of Patent Application Processing and the Office of Finance to the overarching financial realities of this prosecution. Across the 102 nonprovisional patent applications comprising Applicant's master portfolio, the USPTO Office of Finance is currently holding approximately \$35,000.00 to \$42,000.00 USD (with \$32,758.00 USD already empirically verified and credited across 79 examined notices in the Image File Wrappers) in unapplied funds on deposit, without interest and without refund.

Viewing this individual application through an isolated, single-docket lens distorts the financial record by suggesting that Applicant has failed to remit fees. In truth, Applicant has already remitted substantial capital that remains on deposit with the agency. Fundamental administrative equity and MPEP \S{} 509 require that unapplied deposit credits be properly accounted for and applied to maintain procedural timeliness, rather than assessing progressive extension fees that compound to over \$400,000 across the portfolio without prior fee reconciliation.

\subsection*{B. Respectful Request for Allocation of Unapplied Deposit Funds (37 CFR \S{} 1.25(b) \& MPEP \S{} 509)}
Official Office records (specifically the August 27, 2026 Notice of Incomplete Reply and Fee Worksheet WFEE) confirm that the Office of Finance currently holds \textbf{\$416.00 in unapplied funds on deposit} in this application.

Concurrently herewith, Applicant remits the \textbf{\$60.00 processing fee} under 37 CFR \S{} 1.17(i) (Fee Code 2830) via electronic fee payment to satisfy the statutory requirement for the accompanying Request for Suspension under 37 CFR \S{} 1.103(a).

Pursuant to 37 CFR \S{} 1.25(b), MPEP \S{} 509, and MPEP \S{} 710.02(e), Applicant hereby respectfully requests that the Office allocate said \$416.00 unapplied funds on deposit in the following order of procedural priority:
\begin{enumerate}[leftmargin=*]
    \item \textbf{PRIORITY 1 (Timeliness Preservation):} The statutory fee for a Two-Month Extension of Time under 37 CFR \S{} 1.136(a) is \textbf{\$276.00} (Small Entity, Fee Code 2252). Applicant respectfully requests that the Office apply \$276.00 of the available \$416.00 unapplied deposit credit directly to satisfy this extension fee in full as Priority 1. Crediting this amount to the extension fee first establishes that this response is 100\% timely submitted within the extended statutory period under 35 U.S.C. \S{} 133, leaving an unapplied deposit credit balance of \textbf{\$140.00};
    \item \textbf{PRIORITY 2--5 (Pre-Exam Base Fee Deficiency Deferred under 37 CFR \S{} 1.53(f)):} Applicant respectfully requests that the entire remaining deposit credit of \textbf{\$140.00} be applied directly toward the outstanding pre-examination base fees under 37 CFR \S{} 1.16, reducing the net base fee deficiency to \textbf{\$346.00} (\$486.00 balance due less \$140.00). Because this response is timely filed and fully funded as to extensions, Applicant respectfully requests that the Office issue a formal Notice of Fee Deficiency under 37 CFR \S{} 1.53(f) and MPEP \S{} 506 for the remaining base balance, or hold collection in abeyance pursuant to the concurrent Request for Suspension under 37 CFR \S{} 1.103(a).
\end{enumerate}

\section*{IV. CONCLUSION}

All formality requirements having been satisfied, Applicant respectfully requests that the Office of Patent Application Processing (OPAP) advance the application to Technology Center 2100 for examination.

\vspace{0.6cm}
\begin{singlespace}
\noindent Respectfully submitted, \par\vspace{0.8cm}
\underline{\textbf{/Reinaldo Javier Aguiar Marcano/}} \\
\textbf{Reinaldo Javier Aguiar Marcano} \\
Inventor and Pro Se Applicant \\
Mailing Address: 2302 Britton Ridge Drive, Katy, Texas 77494 \\
Telephone: (202) 599-3979 \\
Email: \texttt{rj.aguiar@gmail.com} (USPTO Account) \\
Email: \texttt{rj.aguiar.m@gmail.com} \\
Email: \texttt{reinaldo.aguiar.marcano@isbn9781971481296.com} \\
Date: October 10, 2026
\end{singlespace}

\end{document}
